% ============================================================
% §1 · Introducción — Manual de Usuario K-QGMIP
% ============================================================

\section{Introducción}

\textbf{K-QGMIP} es un \textit{framework} de análisis computacional que permite identificar la
\textbf{Partición de Mínima Información de orden k} (\textit{k-MIP}) en sistemas causales discretos,
en el marco de la \textbf{Teoría de la Información Integrada versión 4.0 (IIT 4.0)}.
El \textit{framework} extiende las estrategias GeoMIP y QNodes —desarrolladas originalmente para
biparticiones ($k = 2$)— al caso general de k-particiones con $k \in \{2, 3, 4, 5\}$,
implementando dos métodos complementarios:

\begin{itemize}
  \item \textbf{KGeoMIP}: estrategia geométrica que aprovecha la representación del espacio de
        estados como un hipercubo $n$-dimensional y reutiliza la tabla de costos de transiciones
        ya calculada para identificar particiones óptimas sin búsqueda exhaustiva.
  \item \textbf{KQNodes}: estrategia submodular que extiende el algoritmo de Queyranne
        (Ordenamiento de Máxima Adyacencia) a k particiones, reduciendo la complejidad
        combinatoria de $\mathcal{O}(S(n,k))$ evaluaciones exhaustivas a una fracción manejable.
\end{itemize}

Desde la perspectiva del usuario, \textbf{K-QGMIP} recibe como entrada la Matriz de Probabilidad
de Transición (TPM) de un sistema, el estado inicial y las máscaras binarias que definen el
subsistema de interés, y devuelve la k-partición óptima junto con su pérdida de información
integrada $\varphi$. El usuario puede seleccionar el método de cálculo (KGeoMIP o KQNodes),
el valor de $k$ y ejecutar el análisis sobre sistemas de distintos tamaños, desde redes de
validación de 10--15 nodos hasta sistemas de escala moderada con 20 o más nodos.

\subsection*{A quién va dirigido este manual}

Este manual está dirigido a \textbf{investigadores, estudiantes y profesionales} que deseen
utilizar K-QGMIP para analizar la estructura causal de sistemas discretos bajo el formalismo
de IIT 4.0. Se asume que el lector posee:

\begin{itemize}
  \item Conocimiento básico de \textbf{Python 3.x} y capacidad de ejecutar scripts desde
        línea de comandos o un entorno de desarrollo integrado (IDE).
  \item Familiaridad general con el concepto de \textbf{matriz de transición probabilística}
        y con la noción de información integrada ($\varphi$), sin necesidad de dominar sus
        fundamentos matemáticos avanzados.
  \item Acceso a una máquina con recursos acordes al tamaño del sistema a analizar
        (ver Sección de Requisitos del Usuario).
\end{itemize}

No se requieren conocimientos previos sobre la teoría de números de Stirling, la
minimización de funciones submodulares ni la geometría de hipercubos; estos conceptos
forman parte del diseño interno del \textit{framework} y el usuario no necesita comprenderlos
para obtener resultados.

\subsection*{Qué no es necesario conocer}

Para utilizar K-QGMIP de forma efectiva, el usuario \textbf{no necesita}:

\begin{itemize}
  \item Conocer el funcionamiento interno de las estrategias KGeoMIP o KQNodes
        (tabla de costos, Ordenamiento de Máxima Adyacencia, NCubos, etc.).
  \item Comprender la demostración de submodularidad de la función de pérdida $\varphi$
        ni el análisis de complejidad asintótica de los algoritmos.
  \item Modificar el código fuente ni configurar las dependencias del entorno de forma manual,
        más allá de los pasos de instalación descritos en la Sección de Instalación y Configuración.
  \item Conocer la arquitectura interna del proyecto (patrones de diseño, jerarquía de clases,
        estructura de módulos) salvo que desee extender o personalizar el \textit{framework}.
\end{itemize}

El presente manual cubre exclusivamente el uso del sistema desde la perspectiva del usuario
final. Para detalles sobre la arquitectura, los algoritmos y el análisis de complejidad,
consulte el \textbf{Manual Técnico} del proyecto.
